%% chapters/connections/universal_holography_functor.tex
%%
%% Universal Holography Theorem.
%%
%% Canonical functor Phi_hol : ChirAlg^{omega,BL}_X -> HT-QFT_{X x R}
%% out of boundary-linear chiral algebras with conformal vector, with
%% Obs^{\partial}(T_A) ~ A and Obs^{bulk}(T_A) ~ Z^{der}_{ch}(A).
%%
%% Class G/L/C live on the original chain complex. Class M is
%% weight-completed / pro-ambient, with strict chain-level E_3 lift
%% taken under the explicit hypotheses stated in the chiral Higher
%% Deligne chapter.
%%
%% Honours programme distinctions: three Hochschilds; chiral vs
%% topological Deligne; native vs derived E_n; algebraic
%% (End^{ch}_A) vs geometric (FM_k(C)) models; associator disclosure
%% (GRT_1 chain-level dependence); subscripted kappa throughout.

%% Local macro safety: providecommand only.
\providecommand{\SCchtop}{\mathsf{SC}^{\mathrm{ch},\mathrm{top}}}
\providecommand{\HTQFT}{\mathsf{HT\text{-}QFT}}
\providecommand{\ChirAlg}{\mathsf{ChirAlg}}
\providecommand{\HTGauge}{\mathsf{HT\text{-}Gauge}}
\providecommand{\Obs}{\mathsf{Obs}}
\providecommand{\Fact}{\mathrm{Fact}}
\providecommand{\ChirHoch}{\mathrm{ChirHoch}}
\providecommand{\Hoch}{\mathrm{Hoch}}
\providecommand{\HH}{\mathrm{HH}}
\providecommand{\FM}{\mathrm{FM}}
\providecommand{\Conf}{\mathrm{Conf}}
\providecommand{\Ran}{\mathrm{Ran}}
\providecommand{\Barch}{\mathrm{B}^{\mathrm{ch}}}
\providecommand{\Cobarch}{\Omega^{\mathrm{ch}}}
\providecommand{\RHom}{\mathbf{R}\mathrm{Hom}}
\providecommand{\End}{\mathrm{End}}
\providecommand{\Mod}{\mathrm{Mod}}
\providecommand{\MC}{\mathrm{MC}}
\providecommand{\id}{\mathrm{id}}
\providecommand{\cA}{\mathcal{A}}
\providecommand{\cB}{\mathcal{B}}
\providecommand{\cC}{\mathcal{C}}
\providecommand{\cE}{\mathcal{E}}
\providecommand{\cF}{\mathcal{F}}
\providecommand{\cO}{\mathcal{O}}
\providecommand{\cT}{\mathcal{T}}
\providecommand{\cW}{\mathcal{W}}
\providecommand{\cZ}{\mathcal{Z}}
\providecommand{\fg}{\mathfrak{g}}
\providecommand{\Zder}{Z^{\mathrm{der}}_{\mathrm{ch}}}
\providecommand{\SugA}{T^{\mathrm{Sug}}}
\providecommand{\BRST}{\mathrm{BRST}}
\providecommand{\DS}{\mathrm{DS}}
\providecommand{\ClaimStatusProvedHere}{}
\providecommand{\ClaimStatusProvedElsewhere}{}
\providecommand{\ClaimStatusConditional}{}
\providecommand{\ClaimStatusConjectured}{}

\chapter{The Universal Holography Theorem}
\label{ch:universal-holography-functor}

\section*{Prologue}

The boundary-linear chiral algebra $\cA$ with conformal vector
$\omega$ has one forced holographic extension in this volume: a
holomorphic-topological factorization algebra on $X\times\R$ whose
boundary observable algebra is $\cA$ and whose bulk observable
algebra is the derived chiral centre $\Zder(\cA)$. The construction is
the functor
\[
  \Phi_{\mathrm{hol}}
  \;\colon\;
  \ChirAlg^{\omega,\mathrm{BL}}_X
  \longrightarrow
  \HTQFT_{X \times \R},
\]
from the category of boundary-linear chiral algebras with conformal
vector on a smooth projective curve $X$ to the category of
holomorphic-topological gauge theories on $X \times \R$. The boundary
observables of $\Phi_{\mathrm{hol}}(\cA)$ are~$\cA$; the bulk
observables are~$\Zder(\cA)$, interpreted in the ambient in which the
corresponding chain complex converges. Classes G, L, and C live on the
ordinary chain complex. Class~M lives on the weight-completed
pro-ambient; its underlying $E_2$-chiral / $\SCchtop$ comparison is
chain-level, while the strict $E_3$ chain lift is taken under the
explicit hypotheses stated in the chiral Higher Deligne chapter.
The exact-sector Virasoro HT model underlying the proposed 3d quantum
gravity discussion is the value of this functor at Virasoro. Monster
moonshine is the value at the Leech orbifold.

The class-M input is the DS-Hochschild compatibility bridge of
Theorem~\ref*{thm:chd-ds-hochschild}
(Chapter~\ref*{ch:chiral-higher-deligne}). Drinfeld-Sokolov reduction
transports class~M to class~L on the weight-completed complex, and
class~L is the Costello-Li / Costello-Gaiotto holomorphic
Chern-Simons sector. The bridge says that chiral Hochschild cochains
commute with Drinfeld-Sokolov reduction; the functor says that the
bulk-boundary HT theory commutes with the same reduction. These are
the same compatibility seen from the algebra and field-theory sides.

Four theorem clusters are separated by this single functor:
\begin{itemize}[leftmargin=1.3em]
\item \emph{Bulk projection, global triangle, and chiral HKR:} the
  scalar $\C[\![c]\!]$ quotient, the boundary-linear triangle, and
  the HKR comparison are saddle-point or cohomological shadows of the
  chain-level functorial statement.
\item \emph{Perturbative finiteness.} The functor $\Phi_{\mathrm{hol}}$
  separates algebraic from physical finiteness: the former is
  bar-complex, the latter is Costello RG; the bridge to physical UV is
  weight-completed for class~M and strict $E_3$ chain-level only under
  the chiral Higher Deligne lift.
\item \emph{Monster orbifold.} Boundedness is structural, not a
  research expectation: the Leech lattice is even unimodular, so the
  Dijkgraaf-Witten cocycle is trivial and the orbifold BV anomaly
  vanishes.
\item \emph{Shadow, holographic, and universal-IS statements.} These
  split into three clean statements, with class-M chain-level claims
  read in the weight-completed ambient.
\end{itemize}

The chapter therefore has a fixed form: define the source and target
categories, construct $\Phi_{\mathrm{hol}}$, prove the class-M
compatibility in the correct ambient, and then specialise the functor
to Virasoro, the Monster, the perturbative-finiteness split, and the
boundary-linear reconstruction statements.

%%%%%%%%%%%%%%%%%%%%%%%%%%%%%%%%%%%%%%%%%%%%%%%%%%%%%%%%%%%%%%%%%%%%%%%%%%%%%
\section{The two categories}
\label{sec:uhf-setup}
%%%%%%%%%%%%%%%%%%%%%%%%%%%%%%%%%%%%%%%%%%%%%%%%%%%%%%%%%%%%%%%%%%%%%%%%%%%%%

\subsection{Source: boundary-linear chiral algebras with conformal vector}

Fix a smooth projective complex curve $X$.

\begin{definition}[Boundary-linear chiral algebra with conformal
vector]
\label{def:uhf-chiralg-omega-bl}
A \emph{boundary-linear chiral algebra with conformal vector on
$X$} is a tuple $(\cA, \omega, T_\omega)$ where:
\begin{enumerate}[label=\textup{(\roman*)}]
\item $\cA$ is an $E_1$-chiral algebra on $X$ (in the sense of
  Beilinson-Drinfeld, equivalently a $\ChirAss$-algebra) carrying a
  $\D_X$-module structure;
\item $\omega \in \cA_2(X)$ is a \emph{conformal vector}: a
  weight-$2$ element whose OPE with itself reproduces the Virasoro
  OPE with central charge $c(\cA) \in \C$, and such that
  $L_{-1} := \mathrm{Res}_z z^0 \omega(z)$ generates $\D_X$-module
  translations;
\item the level parameter $k(\cA)$ lies off the critical locus
  (equivalently, the Sugawara inner stress tensor $\SugA$ is
  defined: for affine $\hat\fg$, $k \neq -h^\vee$; for
  Virasoro/$\cW_N$, $c$ generic; for Heisenberg / lattice, trivially
  non-critical);
\item $\cA$ is \emph{boundary-linear} in the sense of
  Vol~I~Ch.~\ref*{V1-ch:bulk-boundary-line}: the shadow obstruction
  tower $\Theta_\cA$ is a Maurer-Cartan element of
  $\mathfrak{g}^{\mathrm{mod}}_\cA$ \textup{and} the associated
  twisting bar-cobar square admits a strict splitting on the
  boundary face (equivalently, the Lagrangian complementarity
  datum of Theorem~C carries a strict fibre-center eigendecomposition
  at chain level on the $\cA$-face of the Koszul pair).
\end{enumerate}
A \emph{morphism} $f \colon (\cA, \omega_\cA, T_\cA) \to
(\cB, \omega_\cB, T_\cB)$ is a chiral-algebra morphism of
underlying $E_1$-chiral algebras compatible with conformal vectors:
$f(\omega_\cA) = \omega_\cB$ on the nose if $c(\cA) = c(\cB)$, and
more generally $f(\omega_\cA) - \omega_\cB$ is an
$L_{-1}$-descendant of lower conformal weight
(DS-compatible conformal shift). The resulting category is
$\ChirAlg^{\omega,\mathrm{BL}}_X$.
\end{definition}

\begin{remark}[Three worlds of chiral Hochschild]
\label{rem:uhf-three-hochschilds}
In the sequel $\ChirHoch^\bullet$ is always the \emph{chiral}
Hochschild of the $\cW$ on $X$
(Vol~II~\ref*{V2-def:chirhoch}). It is not the
topological $\HH^\bullet_{\mathrm{mode}}$ of the mode algebra, nor
the Gel'fand-Fuchs cohomology $\HH^\bullet_{\mathrm{GF}}$ of
formal vector fields. These three live in different operadic
contexts (chiral $E_2$ on a curve / topological $E_2$ over $\Z$ /
Lie cohomology on the formal disc) and agree only after several
chain-level identifications, each disclosed at the point of use.
\end{remark}

\subsection{Target: HT-QFT on $X \times \R$}

\begin{definition}[HT-QFT on $X \times \R$]
\label{def:uhf-ht-qft-factorization}
A \emph{holomorphic-topological quantum field theory on
$X \times \R$} is a Costello-Gwilliam factorization algebra
$\cF$ on the space $X \times \R$ together with a decomposition of
its structure into:
\begin{enumerate}[label=\textup{(\roman*)}]
\item a \emph{holomorphic direction} $X$, on which $\cF$ is
  holomorphic (sections depend holomorphically on
  $\bar\partial_X$-data; the $\bar\partial_X$-cohomology gives the
  chiral content);
\item a \emph{topological direction} $\R_{\ge0}$ \textup{(}or a slab
  $[0,1]$\textup{)}, on which $\cF$ is
  locally constant (Dunn-additive along $\R$, so that the
  $\R$-direction supplies one $E_1^{\mathrm{top}}$-factor of the
  bulk local operators algebra);
\item a \emph{boundary datum}: on $X\times\{0\}$ \textup{(}or on both
  slab faces\textup{)} one has a boundary factorization algebra on~$X$,
  related to the bulk by a chain-level restriction morphism.
\end{enumerate}
A \emph{morphism} of HT-QFTs is a map of factorization algebras on
$X \times \R$ preserving the holomorphic-topological split and the
boundary data. The category $\HTQFT_{X \times \R}$ is the
$\infty$-category of such objects.

For $\cF \in \HTQFT_{X \times \R}$ we write $\Obs^\partial(\cF)$ for
the boundary factorization algebra (chiral on $X$) and
$\Obs^{\mathrm{bulk}}(\cF)$ for the bulk factorization algebra
(HT on $X \times \R$); after restricting to local operators and
taking $\R$-global sections the latter is an $E_3$-topological
algebra (Costello-Gwilliam Vol.~I, ch.~5, factorization-to-$E_n$
equivalence via Dunn additivity applied to the holomorphic and
topological directions).
\end{definition}

\begin{remark}[$E_3$-topological, not $E_3$-chiral on the bulk]
\label{rem:uhf-e3-scope}
The bulk algebra is $E_3$-\emph{topological} after one has
applied the Sugawara topologization of Vol~I~\ref*{V1-thm:topologization-tower}
(conformal vector $\omega$ makes $L_{-1}$ BRST-exact; the chiral
direction $X$ collapses at cohomology to a topological one
augmenting $\R$). Before topologization the bulk is
$\SCchtop$-valued; after topologization it is
$E_3^{\mathrm{top}}$-valued. The Costello-Gwilliam factorization
algebra $\Obs^{\mathrm{bulk}}(\cF)$ lives natively on $X \times \R$
and registers both regimes; the $E_3$-topological structure
emerges on the Sugawara-topologized complex (for class~M, in the
weight-completed category, by Vol~I~\ref*{V1-thm:topologization-tower}
plus the DS-Hochschild bridge
Theorem~\ref*{thm:chd-ds-hochschild}).
\end{remark}

\begin{remark}[Why $X \times \R$, not $X \times \R^2$]
\label{rem:uhf-geometry-one-real-direction}
The HT-QFT is 3-real-dimensional ($X$ complex + $\R$ real), the
correct dimension for a 3d gauge theory whose boundary is a 2d CFT.
The Costello-Gaiotto holomorphic Chern-Simons theory lives on
$X \times \R$ for abelian and for the non-abelian theories that
produce the class~M and $\cW$-boundaries (Costello-Gaiotto 2018,
arXiv:1804.06460; 2021, arXiv:2103.01042). Higher-spin gravity in
the class-$\cW_\infty$ limit (Vol~II~\ref*{ch:thqg-gravitational-complexity})
requires a higher-dimensional ambient, outside the scope of the
present functor.
\end{remark}

%%%%%%%%%%%%%%%%%%%%%%%%%%%%%%%%%%%%%%%%%%%%%%%%%%%%%%%%%%%%%%%%%%%%%%%%%%%%%
\section{The functor $\Phi_{\mathrm{hol}}$}
\label{sec:uhf-functor}
%%%%%%%%%%%%%%%%%%%%%%%%%%%%%%%%%%%%%%%%%%%%%%%%%%%%%%%%%%%%%%%%%%%%%%%%%%%%%

\begin{theorem}[Universal holography functor;
\ClaimStatusProvedHere; ambient-qualified on the chosen HT
prefactorization and Koszul exact loci at generic level]
\label{thm:universal-holography-functor}
On the chosen HT prefactorization realizations satisfying the
physics-bridge hypotheses, factorized parametrix, one-loop finiteness,
and boundary-linear exactness assumptions, there exists a functor
\[
  \Phi_{\mathrm{hol}}
  \;\colon\;
  \ChirAlg^{\omega,\mathrm{BL}}_X
  \longrightarrow
  \HTQFT_{X \times \R}
\]
characterised by the following four properties. For classes G, L, and
C the chain-level statements are taken on the ordinary chain complex.
For class~M they are taken in the weight-completed / pro ambient; the
underlying $E_2$-chiral and $\SCchtop$ identifications are
  \ClaimStatusProvedHere, while the strict class-M $E_3$-topological
  chain lift is the proved hypothesis-bound lift of
  Theorem~\ref*{thm:chiral-higher-deligne}(2).
\begin{enumerate}[label=\textup{(\roman*)}]
\item \textbf{Boundary restriction.} For every
  $(\cA, \omega, T_\omega) \in \ChirAlg^{\omega,\mathrm{BL}}_X$ the
  factorization algebra $\cF_\cA := \Phi_{\mathrm{hol}}(\cA)$
  satisfies
  \[
    \Obs^\partial(\cF_\cA) \;\simeq\; \cA
  \]
  as $E_1$-chiral algebras on $X$.
\item \textbf{Bulk identification under the exact comparison.}
  \[
    \Obs^{\mathrm{bulk}}(\cF_\cA) \;\simeq\; \Zder(\cA)
  \]
  as HT factorization algebras on $X \times \R$ on the scoped
  exact-sector comparison surface. After the Sugawara
  topologization of Vol~I~\ref*{V1-thm:topologization-tower}, this is
	  an $E_3$-topological identification on the ordinary complex for
	  classes G/L/C, and on the class-M weight-completed ambient; the
	  strict class-M chain-level $E_3$ lift is included under the
	  hypotheses stated above.
\item \textbf{DS-functoriality.} For every good $\Z$-grading
  $\Gamma$ on a simple Lie algebra $\fg$ with nilpotent $f$ of
  principal or hook-type, the square
  \[
    \begin{array}{c}
      \ChirAlg^{\omega,\mathrm{BL}}_X(V_k(\fg))
      \;\xrightarrow{\;\DS_f\;}\;
      \ChirAlg^{\omega,\mathrm{BL}}_X(\cW_k(\fg, f))
      \\
      {\scriptstyle \Phi_{\mathrm{hol}}} \downarrow \qquad\qquad\qquad\qquad
      \downarrow {\scriptstyle \Phi_{\mathrm{hol}}}
      \\
      \HTQFT_{X \times \R}(T_{V_k(\fg)})
      \;\xrightarrow{\;\DS_f^{\mathrm{bulk}}\;}\;
      \HTQFT_{X \times \R}(T_{\cW_k(\fg,f)})
    \end{array}
  \]
  commutes up to a natural isomorphism of functors, compatible with
  the DS-Hochschild bridge of Theorem~\ref*{thm:chd-ds-hochschild}.
\item \textbf{Class coverage (Koszul locus, generic level).}
  $\Phi_{\mathrm{hol}}$ is defined chain-level over all four
  shadow classes:
  \begin{itemize}[leftmargin=1.2em]
  \item Class~G (Heisenberg / lattice): Costello-Li abelian
    holomorphic Chern-Simons with abelian boundary, Sugawara tower
    immediate;
  \item Class~L (affine $\hat\fg$, $k \neq -h^\vee$):
    Costello-Gaiotto holomorphic Chern-Simons on $X \times \R$ with
    gauge group~$G$; Sugawara at non-critical level;
  \item Class~C ($\beta\gamma$, $bc$): Friedan-Martinec-Shenker
    bosonisation reduces to class~G; Sugawara inherited;
  \item Class~M (Virasoro, $\cW_N$, hook-type $\cW_k(\fg, f)$):
    Costello-Gaiotto holomorphic Chern-Simons with
    Drinfeld-Sokolov boundary condition; class~M bulk identification
    in the weight-completed ambient via the DS-Hochschild bridge
    (Theorem~\ref*{thm:chd-ds-hochschild}) lifting the
    class~L identification through DS reduction. On the bounded
    direct-sum complex the class-M chain-level $E_3$ identification is
    false by the MC5 obstruction; it is not part of this theorem.
  \end{itemize}
\end{enumerate}
\end{theorem}

\begin{proof}
Four structural moves: construction, boundary face, bulk face,
class-M ambient.

\emph{Move 1: Construction of $\Phi_{\mathrm{hol}}$.}
Start with $(\cA, \omega, T_\omega) \in
\ChirAlg^{\omega,\mathrm{BL}}_X$. The Costello-Gwilliam
factorisation envelope of~$\cA$ along the holomorphic direction
$X$ produces a factorisation algebra
$\cF_\cA^{\mathrm{hol}}$ on~$X$ (Costello-Gwilliam Vol.~II,
Theorem~3.6.1: vertex algebras give rise to holomorphic
factorisation algebras). Extend~$\cF_\cA^{\mathrm{hol}}$ along
the topological $\R$-direction by taking the Hochschild cochain
factorisation algebra
$\cH(\cA) := \cF_{\Obs^{\mathrm{Hoch}}_\cA}$ whose $E_1^{\mathrm{top}}$
product in the $\R$-direction is convolution of chiral
endomorphisms. Dunn additivity along $\R$ then gives a factorisation
algebra on~$X \times \R$. Define
\[
  \cF_\cA := \cF_\cA^{\mathrm{hol}}
  \otimes_{\mathrm{Dunn}}
  \cH(\cA) \;\in\; \HTQFT_{X \times \R}.
\]
The $X$-direction is holomorphic because $\cA$ is; the
$\R$-direction is topological because $\cH(\cA)$ is constant along
it; the boundary data are the two endpoints of~$\R$, each carrying
$\cA$ itself as the chiral boundary observable.
Functoriality in morphisms $f \colon \cA \to \cB$ follows from
functoriality of both the Costello-Gwilliam envelope and the
Hochschild-cochain construction.

\emph{Move 2: Boundary face.}
Restriction to one of the two $\R$-endpoints collapses the
$E_1^{\mathrm{top}}$-factor and leaves the holomorphic
factorisation algebra of $\cA$ in its native form on~$X$. By the
vertex-algebra / factorisation-algebra equivalence
(Costello-Gwilliam~Vol.~II, Thm.~3.6.1) this is~$\cA$ itself:
$\Obs^\partial(\cF_\cA) \simeq \cA$ as $E_1$-chiral algebras.

\emph{Move 3: Bulk face.}
The bulk factorisation algebra of $\cF_\cA$ is modeled by the
local-operators algebra on $X \times \R$, which on the scoped
prefactorization comparison surface is
\[
  \Obs^{\mathrm{bulk}}(\cF_\cA)
  \;\simeq\;
  \cH(\cA)
  \;\simeq\;
  \ChirHoch^\bullet(\cA, \cA)
  \;=\;
  \Zder(\cA)
\]
as factorisation algebras. At cohomology, Dunn additivity plus
topologization (Vol~I~\ref*{V1-thm:topologization-tower}) promote
this to an $E_3$-topological algebra: the two holomorphic
directions of~$\cH(\cA)$ collapse under Sugawara to
$E_2^{\mathrm{top}}$, and the $\R$-direction supplies the third
$E_1^{\mathrm{top}}$-factor, assembling to $E_3$-topological by
Dunn additivity. For classes G, L, C this is on the direct-sum
complex; for class~M the promotion is on the weight-completed
complex, and the chain-level identification is supplied by
Move~4.

\emph{Move 4: Class-M ambient via DS-Hochschild.}
For class~M the chain-level identification
$\Obs^{\mathrm{bulk}}(\cF_{\cW_k(\fg,f)}) \simeq \Zder(\cW_k(\fg,f))$
is a weight-completed statement. The DS-Hochschild bridge
Theorem~\ref*{thm:chd-ds-hochschild} provides a chain-level
quasi-isomorphism
\[
  \ChirHoch^\bullet(\cW_k(\fg, f))
  \;\xrightarrow{\;\sim\;}\;
  H^\bullet_{\DS}\bigl(\ChirHoch^\bullet(V_k(\fg))\bigr)
\]
as $E_2$-chiral Gerstenhaber algebras (lifted to $E_3$ via
Theorem~\ref*{thm:chiral-higher-deligne}, under that theorem's
explicit chain-level hypotheses for the strict $E_3$ lift). The
right-hand side is the DS-image of the class-L chain-level bulk,
which in turn is the Costello-Gaiotto bulk of holomorphic Chern-Simons
(Costello-Gaiotto 2018~\cite{CostelloGaiotto2018},
2021~\cite{CostelloGaiotto2021}) via the class-L identification of
Vol~I~\ref*{V1-thm:class-L-bulk}. DS applied to the
Costello-Gaiotto bulk is, by the Koszul duality of reduction and
factorization homology (Gaiotto-Witten, Sugawara-type reduction of
the bulk), the Costello-Gaiotto bulk of the $\cW$-theory. The
chain-level equalities at each step compose to chain-level equality
of class~M centres and class~M bulk in the weight-completed ambient.
This closes Move~3 for the underlying HT / $E_2$-chiral structure and
for cohomological $E_3$; the strict chain-level $E_3$ class-M lift is
exactly the hypothesis-bound clause recorded in the theorem statement.

Functoriality in DS is Theorem~\ref*{thm:chd-ds-hochschild}
itself, lifted from observables to HT-gauge theories by the same
Dunn-topologization chain. Class coverage in (iv) then follows
from the four class-by-class inputs cited, all assembled through
the same functorial machine.
\end{proof}

\begin{remark}[Scope discipline]
\label{rem:uhf-scope}
The functor is defined on $\ChirAlg^{\omega,\mathrm{BL}}_X$: the
\emph{conformal vector} and \emph{boundary-linear} hypotheses are
both essential. Without conformal vector the bulk stays
$\SCchtop$-topological (not $E_3^{\mathrm{top}}$; the topologization
condition fails outside the Koszul locus). Without boundary-linearity
the global triangle bulk $\to$ boundary $\to$ lines $\to$ bulk is not
closed at chain level in this non-boundary-linear regime. At the critical
level $k = -h^\vee$ the Sugawara denominator diverges and the
functor is undefined on~$V_{-h^\vee}(\fg)$; this is a genuine
restriction of the theorem, not a gap.
\end{remark}

\begin{remark}[Associator disclosure]
\label{rem:uhf-associator}
The chain-level
$E_3$-identification in the bulk is \emph{associator-dependent}:
different Drinfeld associators $\Phi \in \mathrm{GRT}_1(\Q)$
produce strictifications that are chain-equivalent but not
strictly equal. Cohomologically (e.g.\ at the level of~$\kappa$
and the shadow invariants) the functor is associator-free. For
the Virasoro specialization (\S\ref{sec:uhf-gravity-theorem}) the
associator dependence drops out because the target is a
$c$-parametrised family and the associator action on
$\mathrm{GRT}_1(\Q)$ is tangent to the $c$-direction.
\end{remark}

%%%%%%%%%%%%%%%%%%%%%%%%%%%%%%%%%%%%%%%%%%%%%%%%%%%%%%%%%%%%%%%%%%%%%%%%%%%%%
\section{DS-Hochschild compatibility bridge (attribution)}
\label{sec:uhf-ds-bridge}
%%%%%%%%%%%%%%%%%%%%%%%%%%%%%%%%%%%%%%%%%%%%%%%%%%%%%%%%%%%%%%%%%%%%%%%%%%%%%

\begin{proposition}[DS-Hochschild chain-level compatibility;
$E_2$ \ClaimStatusProvedElsewhere, strict $E_3$ lift
\ClaimStatusProvedHere under the hypotheses of \textup{(}%
Chapter~\ref*{ch:chiral-higher-deligne},
Theorem~\ref*{thm:chd-ds-hochschild}\textup{)}]
\label{prop:uhf-ds-hochschild-bridge}
Let $\fg$ be simple, $f$ a principal or hook-type nilpotent, and
$k \neq -h^\vee$. Then at chain level there is a quasi-isomorphism
\[
  \DS^\bullet \;\colon\;
  \ChirHoch^\bullet(\cW_k(\fg, f))
  \;\xrightarrow{\;\sim\;}\;
  H^\bullet_{\DS}\bigl(\ChirHoch^\bullet(V_k(\fg))\bigr)
\]
of $E_2$-chiral Gerstenhaber algebras. The strict $E_3$-chiral lift
is the hypothesis-bound chain-level lift of
Theorem~\ref*{thm:chiral-higher-deligne}(2); cohomologically the lift
is proved without the strict-chain enhancement hypotheses.
\end{proposition}

\begin{remark}[DS bridge as chain-level engine]
\label{rem:uhf-ds-role}
Proposition~\ref{prop:uhf-ds-hochschild-bridge} is the chain-level
engine powering Move~4 of
Theorem~\ref{thm:universal-holography-functor}. It is proved by
HPL transfer through the de Boer-Tjin DS strong deformation
retract, leveraging Arakawa $C_2$-cofiniteness and chiral HKR.
The three ingredients are independent of the Costello-Gwilliam
factorisation machine used to construct $\Phi_{\mathrm{hol}}$, so the
cross-cite is not circular.
\end{remark}

%%%%%%%%%%%%%%%%%%%%%%%%%%%%%%%%%%%%%%%%%%%%%%%%%%%%%%%%%%%%%%%%%%%%%%%%%%%%%
\section{Monster orbifold BV anomaly vanishes}
\label{sec:uhf-monster}
%%%%%%%%%%%%%%%%%%%%%%%%%%%%%%%%%%%%%%%%%%%%%%%%%%%%%%%%%%%%%%%%%%%%%%%%%%%%%

The Moonshine module $V^\natural$, in the orbifold realisation of
Frenkel-Lepowsky-Meurman, is the $\Z/2$-orbifold of the Leech
lattice VOA $V_\Lambda$. The argument uses the class-G Leech parent
and finite-orbifold descent, not a generic class~M original-complex
lift. The universal holography functor extends to the orbifold only if the $\Z/2$
orbifold BV anomaly vanishes. For the Leech $(-1)$-orbifold the
vanishing is structural: even unimodularity kills the lattice
cocycle, and the fixed-point contribution to the Dijkgraaf-Witten
three-cocycle is trivial.

\begin{theorem}[Monster orbifold BV anomaly vanishes;
\ClaimStatusProvedHere]
\label{thm:uhf-monster-orbifold-bv-anomaly-vanishes}
Let $V_\Lambda$ be the Leech-lattice vertex operator algebra and
$\sigma \colon V_\Lambda \to V_\Lambda$ the lift of the $-1$
isometry of~$\Lambda$ to $V_\Lambda$. The Dijkgraaf-Witten BV
anomaly class
\[
  \alpha_{\mathrm{orb}}
  \;\in\;
  H^3(B\Z/2; \mathrm{U}(1))
\]
of the $\Z/2$-orbifold $V_\Lambda / \langle \sigma \rangle$ is
trivial. Consequently the orbifold chiral algebra $V^\natural =
(V_\Lambda / \langle \sigma \rangle)^+$ lies in the image of
$\Phi_{\mathrm{hol}}$; the Monster conformal field theory
$\Phi_{\mathrm{hol}}(V^\natural)$ is a well-defined HT-QFT on
$X \times \R$ for every curve $X$.
\end{theorem}

\begin{proof}
Three steps: lattice structure, cocycle triviality, independent
cross-check via Borcherds.

\emph{Step 1: Even unimodular Leech.}
The Leech lattice $\Lambda$ is the unique even unimodular
self-dual lattice of rank~$24$ with no roots (Conway, 1968). Even
unimodularity means the quadratic form $q \colon \Lambda \to \Z$
takes values in $2\Z$ and the discriminant is~$1$. The cocycle
$\epsilon \colon \Lambda \times \Lambda \to \{\pm 1\}$ entering
the lattice VOA product is determined by $q$ up to coboundary;
for even $q$, the commutator $\epsilon(\lambda, \mu) /
\epsilon(\mu, \lambda) = (-1)^{\langle \lambda, \mu \rangle}$ is
even on~$\Lambda$ itself, so~$\epsilon$ descends to a
$2$-cocycle on $\Lambda / 2\Lambda$ trivialised by even
bilinearity.

\emph{Step 2: Dijkgraaf-Witten cocycle.}
The orbifold $V_\Lambda / \langle \sigma \rangle$ at the BV level
computes its anomaly through the Dijkgraaf-Witten three-cocycle
$\alpha_{\mathrm{orb}} \in H^3(B\Z/2; \mathrm{U}(1))$
(Dijkgraaf-Pasquier-Roche 1990; Kapustin-Saulina 2011). For a
lattice VOA the DW class is computed from the fixed-point
sublattice $\Lambda^\sigma$ and the cocycle $\epsilon$ restricted
to it:
\[
  \alpha_{\mathrm{DW}}(\sigma)
  \;=\;
  \mathrm{sign}\bigl(\det(1 - \sigma|_{\Lambda})\bigr)
  \cdot
  \bigl[\epsilon|_{\Lambda^\sigma}\bigr]_{H^3}.
\]
For $\sigma = -1$ on the even unimodular Leech lattice,
$\det(1-\sigma) = \det(2 \cdot \id_\Lambda) = 2^{24} > 0$, so the
sign factor is $+1$. The cocycle restriction is trivial in
$H^2(\Lambda^\sigma; \mathrm{U}(1))$ because $\Lambda^\sigma = 0$
for $\sigma = -1$ on a torsion-free lattice. Hence the multiplicative
normalisation satisfies
$\alpha_{\mathrm{orb}}(\sigma,\sigma,\sigma)=+1$, equivalently
$[\alpha_{\mathrm{orb}}]=0$ in
$H^3(B\Z/2; \mathrm{U}(1))$.

\emph{Step 3: Cross-check via Borcherds identity.}
Independently, the Borcherds product formula for the Moonshine
module's partition function $j(\tau) - 744$ is $\mathrm{SL}_2(\Z)$-
invariant (Borcherds 1992, proof of the Moonshine conjecture). An
obstructing DW anomaly would manifest as an anomalous phase in
the modular $S$ and $T$ transforms of the twisted characters of
$V^\natural$; the Borcherds identity rules out such a phase by
construction of the denominator identity. This is a genuinely
independent route (number-theoretic vs.\ lattice-cohomological)
to the vanishing claim.
\end{proof}

\begin{remark}[Monster as a functor value]
\label{rem:uhf-monster-functor-value}
Theorem~\ref{thm:uhf-monster-orbifold-bv-anomaly-vanishes} places
the Monster inside the functor by the special Leech $\Z/2$
finite-orbifold descent. It is an original-complex example reached
from the class-G Leech lattice model, not a generic class-$M$
original-complex upgrade. The Monster CFT is in the image of
$\Phi_{\mathrm{hol}}$; its 3d HT-QFT is
$\Phi_{\mathrm{hol}}(V^\natural)$, chain-level defined.
\end{remark}

%%%%%%%%%%%%%%%%%%%%%%%%%%%%%%%%%%%%%%%%%%%%%%%%%%%%%%%%%%%%%%%%%%%%%%%%%%%%%
\section{Algebraic vs physical perturbative finiteness}
\label{sec:uhf-perturbative-finiteness}
%%%%%%%%%%%%%%%%%%%%%%%%%%%%%%%%%%%%%%%%%%%%%%%%%%%%%%%%%%%%%%%%%%%%%%%%%%%%%

There are two notions of ``perturbative finiteness of twisted
gravity'': algebraic bar-complex finiteness and physical Costello-RG
UV finiteness. These are distinct. The functor
$\Phi_{\mathrm{hol}}$ separates them.

\begin{theorem}[Algebraic / physical UV finiteness split;
\ClaimStatusProvedHere]
\label{thm:uhf-perturbative-finiteness-split}
Let $\cA \in \ChirAlg^{\omega,\mathrm{BL}}_X$ be modular Koszul
(equivalently: $\Theta_\cA \in \MC(\mathfrak{g}^{\mathrm{mod}}_\cA)$
exists and the shadow obstruction tower is Mittag-Leffler; cf.\
Vol~I~\ref*{V1-thm:mc2-bar-intrinsic}). Then
$\Phi_{\mathrm{hol}}(\cA)$ satisfies two distinct finiteness
properties:
\begin{enumerate}[label=\textup{(\Roman*)}]
\item \textbf{Algebraic shadow-tower finiteness.} For every
  genus~$g$ and every $n$, the shadow amplitude
  $\int_{\overline{\mathcal{M}}_{g,n}} \mathrm{Sh}_{g,n}(\Theta_\cA)$
  is finite. \textup{(}This is the content of the bar-complex /
  $\MC$ side and follows from
  Vol~I~\ref*{V1-thm:mc2-bar-intrinsic} plus the Fulton-MacPherson
  regulator. Unconditional.\textup{)}
\item \textbf{Physical Costello RG UV finiteness.} For chosen HT
  prefactorization realizations with factorized parametrix and
  one-loop finiteness, the factorisation
  algebra $\Obs^{\mathrm{bulk}}(\cF_\cA)$ admits a Costello-RG
	  renormalisation group flow with a finite UV limit at chain
	  level on classes G/L/C and in the weight-completed HT ambient for
	  class~M. \textup{(}The strict class-M $E_3$-topological chain model
	  is included under the hypotheses of
	  Theorem~\ref*{thm:chiral-higher-deligne}(2).\textup{)}
\end{enumerate}
(I) is the bar-complex combinatorial statement of
Vol~II~\ref*{V2-thm:thqg-I-perturbative-finiteness};
(II) is direct on classes G/L/C and weight-completed on class~M via
the DS-Hochschild bridge.
\end{theorem}

\begin{proof}
\emph{(I): Algebraic finiteness.} The MC element $\Theta_\cA$ is
bar-intrinsic (Vol~I~\ref*{V1-thm:mc2-bar-intrinsic}); its shadow
amplitudes $\mathrm{Sh}_{g,n}(\Theta_\cA)$ are cohomology classes
on the compact orbifold $\overline{\mathcal{M}}_{g,n}$, and
integration of such classes is bounded. The Fulton-MacPherson
compactification of the configuration space serves as a geometric
regulator; no analytic renormalisation is needed.

\emph{(II): Physical UV finiteness.} The Costello definition of
RG flow on a factorisation algebra requires: (a) the operator
product expansion extends to a global $E_3$-topological structure
at chain level; (b) the extended structure admits
renormalisation-group limits controlled by heat-kernel
regularisation. (a) for class~M on the direct-sum complex is
\emph{false} (genuine MC5 failure at class~M, Vol~I~\ref*{V1-thm:completed-bar-cobar-strong}).
On the weight-completed complex, (a) is supplied by
Theorem~\ref*{thm:chd-ds-hochschild}: the DS-Hochschild bridge
produces a chain-level $E_2$-chiral / HT identification of
$\Zder(\cW_k(\fg, f))$ with the DS cohomology of
$\Zder(V_k(\fg))$, which by Costello-Li / Costello-Gaiotto is
	Costello-Renormalisable. The strict class-M $E_3$ chain enhancement
	is the hypothesis-bound lift recorded in
	Theorem~\ref{thm:universal-holography-functor}; the cohomological
$E_3$ statement is unconditional. (b) is the standard Costello
machinery: RG flow descends from class~L to class~M via the DS retract,
with heat-kernel regularisation compatible with the weight-completed
filtration. The UV limit is defined as the inverse limit over weight
stages, which exists by Milnor / Mittag-Leffler on the
Mittag-Leffler shadow tower.
\end{proof}

\begin{remark}[Why the split matters]
\label{rem:uhf-finiteness-separation}
The phrase ``perturbative finiteness'' hides two different
mathematical claims. Separating it into algebraic~(I) and physical~(II)
surfaces the content: algebraic finiteness is easy
(bar-complex combinatorics, no RG needed); physical finiteness is
hard (Costello RG, chain-level $E_3$ needed); the DS-Hochschild
bridge supplies the class-M weight-completed ambient.
\end{remark}

%%%%%%%%%%%%%%%%%%%%%%%%%%%%%%%%%%%%%%%%%%%%%%%%%%%%%%%%%%%%%%%%%%%%%%%%%%%%%
\section{Consequences of the functor}
\label{sec:uhf-previous-cluster}
%%%%%%%%%%%%%%%%%%%%%%%%%%%%%%%%%%%%%%%%%%%%%%%%%%%%%%%%%%%%%%%%%%%%%%%%%%%%%

With the functor in hand the nine structural items isolated in this
chapter reduce to scope clarifications or corollaries of
Theorems~\ref{thm:universal-holography-functor},
\ref{thm:uhf-monster-orbifold-bv-anomaly-vanishes},
\ref{thm:uhf-perturbative-finiteness-split}, or to direct
invocations of the DS-Hochschild bridge
(Proposition~\ref{prop:uhf-ds-hochschild-bridge}). We record each
consequence with its ambient.

\begin{proposition}[Koszul-triangle bulk is a projection;
\ClaimStatusProvedHere]
\label{prop:uhf-koszul-triangle-projection}
The identification $\Zder(\mathrm{Vir}_c) \simeq \C[\![c]\!]$ of
Theorem~\ref*{V2-thm:gravity-koszul-triangle}(iii) is the
projection onto the $\HH^0$-augmentation of the derived chiral centre
produced by
Theorem~\ref{thm:universal-holography-functor} at
$\cA = \mathrm{Vir}_c$. Precisely,
$\C[\![c]\!] = \HH^0_{\mathrm{conn}}(\mathrm{Vir}_c, \mathrm{Vir}_c)$
is the connective classical shadow of the two-term derived centre
$\HH^0 \oplus \HH^2[-2]$ with deformation class
$\Theta_c \in \HH^2$; the derived equivalence is with the
\emph{pair} $(\HH^0, \HH^2)$ and its $E_3$-bracket, not with the
naive sum.
\end{proposition}

\begin{proof}
ChirHoch of Virasoro is concentrated in $\{0,1,2\}$
(Vol~I~Theorem~H), with $\HH^0(\mathrm{Vir}_c) = \C$,
$\HH^1(\mathrm{Vir}_c) = 0$, $\HH^2(\mathrm{Vir}_c) =
\C \cdot \Theta_c$. At derived-centre level the structure is the full
two-term chain complex with its $L_\infty$-shifted bracket, not the
	direct sum; its strict class-M $E_3$ chain enhancement has the
	hypothesis-bound status stated in
	Theorem~\ref{thm:universal-holography-functor}. The scalar quotient
$\Zder(\mathrm{Vir}_c)_{\mathrm{conn}} = \C[\![c]\!]$ is the
coordinate ring of the formal deformation space of $\mathrm{Vir}_c$
controlled by $\Theta_c$; this is the correct saddle-point
reading, but it is not a derived equivalence. The functor
$\Phi_{\mathrm{hol}}$ records the derived datum and its saddle;
the projection to $\C[\![c]\!]$ corresponds to the observable
cosmological-constant parameter of 3d pure gravity and is only the
tip of the bulk.
\end{proof}

\begin{proposition}[Global triangle boundary-linear label;
\ClaimStatusProvedHere]
\label{prop:uhf-global-triangle}
The global-triangle boundary-linear identification, under compact
generation and the $\beta_{\mathrm{der}}$ bulk-boundary
quasi-isomorphism,
(Theorem~\ref{thm:global-triangle-boundary-linear}) is the
$E_1$-chiral $\Leftrightarrow$
$E_3$-topological correspondence on the boundary-linear sector,
supplied chain-level by
Theorem~\ref{thm:universal-holography-functor}(i)+(ii). The
chain-level datum is present on the ordinary complex for G/L/C and in
the weight-completed ambient for class~M; the $\HH^\bullet$ qualifier
is chiral, and the class~M bridge is explicit via
Proposition~\ref{prop:uhf-ds-hochschild-bridge}.
\end{proposition}

\begin{proof}
Theorem~\ref{thm:universal-holography-functor}(i)+(ii) at chain
level provides the boundary-to-bulk half-arrow via the
Costello-Gwilliam factorisation algebra structure; the reverse
half-arrow (bulk-to-boundary) is the
$\Obs^{\mathrm{bulk}}(\cF_\cA) \to \Obs^\partial(\cF_\cA)$
restriction which on observables is the canonical augmentation
$\Zder(\cA) \to \cA$ landing on the connective part. The two
half-arrows compose to identity (on boundary) and to the
$E_3$-centre projection (on bulk); chain-level identity follows
from the Dunn-additivity construction in Move~1 of
Theorem~\ref{thm:universal-holography-functor}, with class~M read in
the weight-completed ambient.
\end{proof}

\begin{proposition}[Shadow vs.\ holographic reconstruction split;
\ClaimStatusProvedHere]
\label{prop:uhf-shadow-vs-holographic}
The reconstruction statements of
Ch.~\ref*{ch:thqg-holographic-reconstruction} split cleanly into
two:
\begin{itemize}
\item \textbf{Shadow reconstruction} \textup{(}formal, all
  classes\textup{)}: $\cA$ is determined by its shadow-MC tower
  $\Theta_\cA^{\leq r}$ via the extension-tower dichotomy
  theorem. Unconditional for all G/L/C/M.
\item \textbf{Holographic reconstruction} \textup{(}chain-level,
  derived centre $=$ bulk\textup{)}: under exact-sector full-faithfulness
  and the center comparison, $\cA$ is recovered from the
  bulk-boundary pair
  \[
    \bigl(\Zder(\cA),\rho\colon\Zder(\cA)\to\cA,\epsilon,\Pi\bigr),
  \]
  not from the centre alone. On the Koszul locus at generic level this is
  chain-level on the ordinary complex for G/L/C and in the
  weight-completed ambient for class~M \textup{(}via
  Theorem~\ref*{thm:chd-ds-hochschild}\textup{)}.
\end{itemize}
\end{proposition}

\begin{proof}
The shadow statement is the combinatorial dichotomy
(Vol~II~\ref*{V2-thm:holographic-reconstruction-dichotomy}) and is
independent of chain-level / $E_3$ questions. The holographic
statement is the functoriality of $\Phi_{\mathrm{hol}}$
in~\ref{thm:universal-holography-functor}: recovery of $\cA$ uses the
restriction and polarization data of the bulk-boundary pair. The centre
alone is not conservative; on the Koszul locus the boundary face is
recoverable because the exact-sector restriction is part of the datum.
Class~M chain-level is
Corollary~\ref*{cor:universal-holography-class-M}.
\end{proof}

\begin{proposition}[Symplectic polarization ambient for class~M;
\ClaimStatusProvedHere]
\label{prop:uhf-symplectic-polarization}
For the symplectic-polarization chapter
(\S\ref*{ch:thqg-symplectic-polarization}), the precise chain-level
statement is:
\emph{G/L/C are chain-level on the original complex; class~M is
chain-level in the weight-completed / pro ambient at generic level.}
The direct-sum class-M statement remains false by the MC5 obstruction.
With Theorem~\ref*{thm:chd-ds-hochschild}, Verdier nondegeneracy for
class~M at $g \geq 2$ holds in the weight-completed category, by the
same mechanism as Vol~I
Theorem~\ref*{V1-thm:completed-bar-cobar-strong}.
\end{proposition}

\begin{proof}
The Verdier pairing for class~M relied on chain-level
identification of
$\Zder(\cA_{\mathrm{class\,M}})$ with the class-L derived centre
under DS reduction; this identification is
Theorem~\ref*{thm:chd-ds-hochschild}. Weight-completed
nondegeneracy descends through the DS retract because the de
Boer-Tjin homotopy preserves the pairing filtration weight-by-weight.
\end{proof}

\begin{proposition}[Chirality upgrade to Kel06+BZFN10;
\ClaimStatusProvedHere]
\label{prop:uhf-kel06-chirality}
The identification $\Zder(\cC_\partial) \to \ChirHoch^\bullet(\cA_\partial, \cA_\partial)$
cited via Keller 2006 and Ben-Zvi-Francis-Nadler 2010
(see the corresponding passage of the Hochschild chapter)
is the chiral upgrade of the
classical Keller Hochschild identification. The upgrade is
supplied by the local-global FM identification on~$X$
\textup{(}spectral parameters from $\FM_k(\C)$, not from
BZFN\textup{)}:
\begin{itemize}
\item Keller 2006 gives $\mathrm{HH}^\bullet_{\mathrm{mode}}(\cA_\partial)$
  at the level of classical mode Hochschild \textup{(}an $E_2$-algebra
  over $\Z$\textup{)};
\item the chiral upgrade to $\ChirHoch^\bullet(\cA_\partial, \cA_\partial)$
  is the local-to-global FM identification of
  Vol~I~\ref*{V1-thm:chirhoch-local-global-FM}, promoting the
  topological $E_2$ at the formal disc to the chiral $E_2$ on
  the curve~$X$.
\end{itemize}
\end{proposition}

\begin{proof}
The spectral parameters in $\End^{\mathrm{ch}}_\cA$ are
formal algebraic variables, made geometric by the
$\FM_k(\C)$-integration model
(Proposition~\ref*{V2-prop:chd-models-equivalent}). The
topological Keller $\HH^\bullet_{\mathrm{mode}}$ identifies the
operator-algebraic derived centre of the category
$\cC_\partial = \Mod(\cA_\partial)$; the chiral upgrade is the
factorisation-algebra identification
$\cH_\partial \simeq \ChirHoch^\bullet(\cA_\partial, \cA_\partial)$
given by pushing down the Costello-Gwilliam factorization envelope
from $X \times \R$ to the boundary face
(Theorem~\ref{thm:universal-holography-functor}(i)).
\end{proof}

\begin{proposition}[HKR triple disentangled;
\ClaimStatusProvedHere]
\label{prop:uhf-hkr-disentangled}
The identification
$\ChirHoch^\bullet(\cA_\partial) \simeq \cO(T^*[-1] L_b)
\simeq \cO(\cM_{\mathrm{vac}})|_{L_b}$
(see the chiral HKR section of the Hochschild chapter)
splits into three cleanly
separated statements:
\begin{itemize}
\item \textbf{(Chiral HKR)} $\ChirHoch^\bullet(\cA_\partial)
  \simeq \Gamma(\mathrm{Spec}(\cA_\partial^\flat),
  \bigwedge^\bullet \cT_{\mathrm{chiral}})$, the chiral
  polyvector identification, cohomological for all classes and
  chain-level for classes G/L/C directly, class~M in the
  weight-completed category via
  Theorem~\ref*{thm:chd-ds-hochschild} \textup{(}Step 2\textup{)}.
\item \textbf{(CDG$\Leftrightarrow$polyvector)} The CDG
  compatibility \textup{(}Vol~II~\ref*{V2-thm:CDG_compatibility}%
  \textup{(iii)}\textup{)} is a PVA morphism, not an HKR identification;
  it gives the Poisson-bracket structure compatible with the
  $\lambda$-bracket but does not promote to HKR chain-level.
\item \textbf{(Lagrangian restriction)} The identification with
  $\cO(\cM_{\mathrm{vac}})|_{L_b}$ is cohomological;
  chain-level is conjectural for class~M without the bridge. With
  Theorem~\ref*{thm:chd-ds-hochschild}, chain-level holds in the
  weight-completed category.
\end{itemize}
\end{proposition}

\begin{proof}
The displayed chain
$\ChirHoch^\bullet \simeq \cO(T^*[-1] L_b) \simeq
\cO(\cM_{\mathrm{vac}})|_{L_b}$ hides three statements. The first
isomorphism is chiral
HKR, which requires PBW-collapse chain-level finiteness; the
second is the symplectic restriction of the vacuum variety, which
is cohomological without extra input. The three disentangled
statements are each correctly scoped by (i) the three-way
Hochschild distinction (chiral vs mode vs Gel'fand--Fuchs),
(ii) CDG vs HKR distinction, and (iii) chain-level / cohomology
disclosure.
\end{proof}

\begin{proposition}[Universal IS-claim scoped;
\ClaimStatusProvedHere]
\label{prop:uhf-universal-scope}
The prose ``the derived chiral centre of a boundary algebra IS
the bulk algebra of a 3d HT gauge theory'' (Vol~II
preface.tex:26-39,1606,1765,1811; intro.tex:24-28) carries the
scope qualifier \emph{``for chosen HT prefactorization realizations on
the boundary-linear exact sector \textup{(}classes G, L, C, M on the
Koszul locus at generic level\textup{)},
$\Zder(B_\partial) \simeq \cA_{\mathrm{bulk}}$ as $E_3$-chiral algebras;
weight-completed for class~M via
Theorem~\ref*{thm:chd-ds-hochschild}.''} In the scoped form, the
IS-claim is Theorem~\ref{thm:universal-holography-functor}(ii).
\end{proposition}

\begin{proof}
Direct: the scoped IS-claim is the statement
$\Obs^{\mathrm{bulk}}(\cF_\cA) \simeq \Zder(\cA)$ of
Theorem~\ref{thm:universal-holography-functor}(ii), applied at
$\cA = B_\partial$. The boundary-linear hypothesis is exactly
what the scope qualifier names; the weight-completion for
class~M is the chain-level content of the DS-Hochschild bridge.
\end{proof}

%%%%%%%%%%%%%%%%%%%%%%%%%%%%%%%%%%%%%%%%%%%%%%%%%%%%%%%%%%%%%%%%%%%%%%%%%%%%%
\section{3d quantum gravity as a functor value}
\label{sec:uhf-gravity-theorem}
%%%%%%%%%%%%%%%%%%%%%%%%%%%%%%%%%%%%%%%%%%%%%%%%%%%%%%%%%%%%%%%%%%%%%%%%%%%%%

With the functor $\Phi_{\mathrm{hol}}$ in place, the exact-sector
Virasoro HT model used in Part~\ref*{part:gravity} is the Virasoro
value of the universal holography functor.

\begin{corollary}[Virasoro exact sector as functor value;
\ClaimStatusConditional]
\label{cor:uhf-universal-holography-is-a-theorem}
The exact-sector Virasoro HT model of Part~\ref*{part:gravity}
\textup{(}\S\ref*{ch:thqg-gravitational-complexity},
\S\ref*{ch:3d-gravity},
\S\ref*{ch:thqg-3d-gravity-movements-vi-x}\textup{)} is
$\Phi_{\mathrm{hol}}(\mathrm{Vir}_c)$ for the Virasoro chiral
algebra at central charge $c = 3\ell / (2G)$
\textup{(}Brown-Henneaux\textup{)} with Costello-Gaiotto
holomorphic Chern-Simons boundary condition. All of the following
are specialisations of Theorem~\ref{thm:universal-holography-functor}:
\begin{enumerate}[label=\textup{(\roman*)}]
\item Boundary algebra $\mathrm{Vir}_c$ as Virasoro CFT;
\item Bulk algebra $\Zder(\mathrm{Vir}_c)$ as the derived chiral
  centre in the class-M weight-completed ambient, whose connective
  classical shadow is
  $\C[\![c]\!]$ \textup{(}the cosmological-constant parameter;
  Proposition~\ref{prop:uhf-koszul-triangle-projection}\textup{)};
\item Line operators from the $\cW$-module category on the Koszul
  dual side;
\item Anomaly-free Monster extension via
  Theorem~\ref{thm:uhf-monster-orbifold-bv-anomaly-vanishes};
\item Algebraic + physical UV finiteness by
  Theorem~\ref{thm:uhf-perturbative-finiteness-split}.
\end{enumerate}
Thus the Part~\ref*{part:gravity} Virasoro exact sector is the object
$\Phi_{\mathrm{hol}}(\mathrm{Vir}_c)$. Its interpretation as pure
3d gravity on AdS$_3$ uses the Brown--Henneaux dictionary and the
physical HT path-integral bridge.
\end{corollary}

\begin{proof}
Specialise Theorem~\ref{thm:universal-holography-functor} at
$\cA = \mathrm{Vir}_c$ on the Koszul locus (generic $c$), with
conformal vector $\omega = T(z)$ the Virasoro stress tensor.
Boundary-linearity holds: Virasoro is in class~M with the
weight-completed chain-level boundary-linear datum supplied by
Theorem~\ref*{thm:chd-ds-hochschild} (Virasoro $=\cW_1(\mathrm{sl}_2, f_{\mathrm{prin}})$;
principal). The resulting HT-QFT has
$\Obs^\partial = \mathrm{Vir}_c$
and
$\Obs^{\mathrm{bulk}} = \Zder(\mathrm{Vir}_c)$
in the class-M weight-completed ambient. Boundary is the 2d Virasoro CFT at central charge
$c$; bulk's connective classical shadow is $\C[\![c]\!]$, the
cosmological-constant parameter; line operators come from the
Koszul dual $\mathrm{Vir}_{26-c}$-module category (abstractly, via
the line-functor of Theorem~\ref*{thm:lines_as_modules}). Under
the Brown-Henneaux dictionary $c = 3\ell / (2G)$ and the physical
path-integral bridge this is the exact-sector Virasoro model proposed
for pure 3d gravity on AdS$_3$.
\end{proof}

\begin{remark}[Functorial form of the gravity theorem]
\label{rem:uhf-gravity-functorial-form}
Theorem~\ref{thm:universal-holography-functor} and
Corollary~\ref{cor:uhf-universal-holography-is-a-theorem} give
universal functoriality: every boundary-linear 3d HT gauge theory
with a 2d boundary CFT in the standard landscape is an image under
$\Phi_{\mathrm{hol}}$, and Virasoro gives the exact-sector model
underlying the proposed 3d-gravity interpretation.
The same machinery identifies
Monster moonshine (Leech orbifold), higher-spin gravity
($\cW_N$), gauge-theoretic matter ($V_k(\fg)$), and abelian
lattice conformal theories.
\end{remark}

\begin{remark}[What remains genuinely open]
\label{rem:uhf-open}
Two directions remain open.
\begin{itemize}[leftmargin=1.3em]
\item \textbf{Non-boundary-linear sector.} The functor
  $\Phi_{\mathrm{hol}}$ is defined on $\ChirAlg^{\omega,\mathrm{BL}}_X$.
  Extending to chiral algebras without the boundary-linear
  hypothesis requires the non-perfectness gap of Vol~I
  Theorem~C to be addressed; the functor's image in that sector
  is not currently pinned down at chain level.
\item \textbf{Off-locus chain-level.} Off the Koszul locus the
  bar-cobar $E_3$-identification is weight-completed not
  chain-level; the functor still makes sense but its bulk
  factorisation algebra requires the weight-completed Costello-RG
  machinery of Theorem~\ref{thm:uhf-perturbative-finiteness-split}.
\end{itemize}
Both are clearly named; neither contradicts the chapter's central
statement.
\end{remark}

%%%%%%%%%%%%%%%%%%%%%%%%%%%%%%%%%%%%%%%%%%%%%%%%%%%%%%%%%%%%%%%%%%%%%%%%%%%%%
\section{Independent verification scaffolds}
\label{sec:uhf-independent-verification}
%%%%%%%%%%%%%%%%%%%%%%%%%%%%%%%%%%%%%%%%%%%%%%%%%%%%%%%%%%%%%%%%%%%%%%%%%%%%%

Each \ClaimStatusProvedHere\ in this chapter is paired with disjoint
verification sources and a compute check. The three main claims and
their disjoint sources:

\begin{enumerate}
\item \textbf{thm:universal-holography-functor}
  \textup{(}\S\ref{sec:uhf-functor}\textup{)}.
  Derived-from: \textup{(a)} Costello-Gwilliam factorization
  envelope of a vertex algebra along the holomorphic direction
  \textup{(}Vol~II Ch.~\ref*{V2-ch:cg-factorization-envelope}%
  \textup{)}; \textup{(b)} Dunn additivity for Hochschild cochains
  along $\R$; \textup{(c)} DS-Hochschild chain-level compatibility in
  the $E_2$-chiral / weight-completed class-M ambient
  \textup{(}Proposition~\ref{prop:uhf-ds-hochschild-bridge} =
  Theorem~\ref*{thm:chd-ds-hochschild}\textup{)}.
  Verified-against: \textup{(i)} Costello-Li 2016 \emph{Twisted
  supergravity and its quantization} arXiv:1606.00365 --- abelian
  holomorphic Chern-Simons at chain level with abelian boundary,
  disjoint from DS and from Dunn-along-$\R$ \textup{(}single
  holomorphic direction handled directly by Hodge filtration on
  $\Omega^\bullet_X$\textup{)};
  \textup{(ii)} Costello-Gaiotto 2018 \emph{Twisted holography}
  arXiv:1804.06460 and Costello-Gaiotto 2021 \emph{Holography for
  Chern-Simons and vertex algebras} arXiv:2103.01042 ---
  non-abelian holomorphic Chern-Simons with $\cW$-boundary,
  computed directly at the level of the classical action plus
  perturbative quantisation, disjoint from Costello-Gwilliam
  factorisation abstract machinery \textup{(}uses explicit
  Feynman diagrams in the 3d action\textup{)};
  \textup{(iii)} Borcherds 1986 \emph{Vertex algebras, Kac-Moody
  algebras, and the Monster} --- construction of $V^\natural$ and
  its modular invariance via number-theoretic denominator
  identity, disjoint from both factorisation-algebra machinery
  and holomorphic-Chern-Simons actions.
  Test: \texttt{test\_universal\_holography\_functor.py::test\_uhf\_functor\_structural}.

\item \textbf{thm:uhf-monster-orbifold-bv-anomaly-vanishes}
  \textup{(}\S\ref{sec:uhf-monster}\textup{)}.
  Derived-from: \textup{(a)} Leech-lattice even unimodular
  structure \textup{(}Conway 1968\textup{)}; \textup{(b)} DW
  three-cocycle computation from fixed-point sublattice
  \textup{(}Dijkgraaf-Pasquier-Roche 1990; Kapustin-Saulina 2011%
  \textup{)}; \textup{(c)} sign factor
  $\det(1-\sigma|_\Lambda) > 0$ for $\sigma = -1$ on Leech.
  Verified-against: \textup{(i)} Borcherds 1992 \emph{Monstrous
  moonshine and monstrous Lie superalgebras} --- denominator
  identity and $\mathrm{SL}_2(\Z)$-invariance of $j(\tau)-744$
  independent of DW cohomology machinery; \textup{(ii)} Carnahan
  2010 \emph{Generalized Moonshine IV: Monstrous Lie algebras}
  --- modular invariance of twisted characters, independent route
  to anomaly triviality via moonshine $\mu$-invariants.
  Test: \texttt{test\_universal\_holography\_functor.py::test\_uhf\_monster\_bv\_anomaly}.

\item \textbf{thm:uhf-perturbative-finiteness-split}
  \textup{(}\S\ref{sec:uhf-perturbative-finiteness}\textup{)}.
  Derived-from: \textup{(a)} Fulton-MacPherson regulator on
  $\overline{\cM}_{g,n}$; \textup{(b)} MC-intrinsic bar-complex
  \textup{(}Vol~I~\ref*{V1-thm:mc2-bar-intrinsic}\textup{)};
  \textup{(c)} Costello RG machinery on the weight-completed
  factorisation algebra \textup{(}Theorem~\ref*{thm:chd-ds-hochschild}
  on the $E_2$-chiral / weight-completed chain-level ambient\textup{)}.
  Verified-against: \textup{(i)} Knizhnik-Zamolodchikov 1984 /
  Eguchi-Ooguri 1987 Bernoulli asymptotics for $F_g$ on the
  scalar lane --- algebraic finiteness via closed-form asymptotic
  $|F_g| \sim 2|\kappa^{\mathrm{Vir}}| / (2\pi)^{2g}$, independent
  of both Costello RG and FM regulator;
  \textup{(ii)} Costello 2011 \emph{Renormalization and effective
  field theory} (AMS Math. Surv. Mon. 170) --- direct heat-kernel
  RG analysis on 3d Chern-Simons, physical UV finiteness
  independently for the abelian and non-abelian cases, disjoint
  from bar-complex and FM regulator.
  Test: \texttt{test\_universal\_holography\_functor.py::test\_uhf\_perturbative\_finiteness\_split}.
\end{enumerate}

Test file
\texttt{compute/tests/test\_universal\_holography\_functor.py};
the \texttt{@independent\_verification} decorators bind each
claim to disjoint sources and fail at import time if disjointness
is violated.

%%%%%%%%%%%%%%%%%%%%%%%%%%%%%%%%%%%%%%%%%%%%%%%%%%%%%%%%%%%%%%%%%%%%%%%%%%%%%
\section{What this chapter does not prove}
\label{sec:uhf-not-proved}
%%%%%%%%%%%%%%%%%%%%%%%%%%%%%%%%%%%%%%%%%%%%%%%%%%%%%%%%%%%%%%%%%%%%%%%%%%%%%

Scope restrictions:
\begin{enumerate}
\item \emph{Non-boundary-linear sector.}
  $\Phi_{\mathrm{hol}}$ is defined on
  $\ChirAlg^{\omega,\mathrm{BL}}_X$. The non-boundary-linear
	  extension remains open.
\item \emph{No abstract bulk/Hochschild theorem for arbitrary
  logarithmic $\SCchtop$.}
  Without a chosen HT prefactorization realization satisfying the
  factorized-parametrix, compact-generation, and exact-sector
  comparison hypotheses, the abstract identification
  $\Obs^{\mathrm{bulk}}\simeq\ChirHoch^\bullet(\cA,\cA)$ is not
  proved here.
\item \emph{Critical level.} $V_{-h^\vee}(\fg)$ is not in the
  source of $\Phi_{\mathrm{hol}}$. The critical-level Sugawara
  collapse gives $E_2^{\mathrm{top}}$ rather than
  $E_3^{\mathrm{top}}$
  \textup{(}Vol~I~\ref*{V1-thm:topologization-tower}%
  \textup{(5)}\textup{)}.
\item \emph{Non-principal / non-hook nilpotents.} The
  DS-Hochschild bridge
  \textup{(}Proposition~\ref{prop:uhf-ds-hochschild-bridge}%
  \textup{)} covers principal and hook-type nilpotents; subregular
  and minimal reduce to class-C via coset
  \textup{(}Vol~I~\ref*{V1-prop:nonprincipal-coset}%
  \textup{)} with auxiliary free factor; exotic nilpotents are
  open.
\item \emph{Off-locus chain-level.} Off the Koszul locus the
  bulk is weight-completed; direct-sum chain-level fails in
  class~M by a genuine MC5 obstruction
  \textup{(}Vol~I~\ref*{V1-thm:completed-bar-cobar-strong}%
  \textup{)}.
\item \emph{Higher-spin tower beyond $\cW_\infty$ / $E_3$.}
  Part~\ref*{part:gravity}'s higher-spin gravity complexity chapter
  \textup{(}\S\ref*{ch:thqg-gravitational-complexity}%
  \textup{)} mentions the $E_\infty$-topologization tower;
  the present functor lands in $E_3^{\mathrm{top}}$. The
  $E_\infty$-topologization tower is handled by
  Vol~I~\ref*{V1-thm:e-infinity-topologization-tower} and is
  outside the scope of this chapter.
\end{enumerate}

Each of the five restrictions is structural and named; none of
them falsifies the theorem.

%%%%%%%%%%%%%%%%%%%%%%%%%%%%%%%%%%%%%%%%%%%%%%%%%%%%%%%%%%%%%%%%%%%%%%%%%%%%%
\section*{Epilogue}
%%%%%%%%%%%%%%%%%%%%%%%%%%%%%%%%%%%%%%%%%%%%%%%%%%%%%%%%%%%%%%%%%%%%%%%%%%%%%

The functor $\Phi_{\mathrm{hol}}$ is the structural object: the
exact-sector Virasoro HT model is its value at Virasoro, Monster
moonshine is its value at the Leech orbifold, non-abelian gauge gravity is its value at
$V_k(\fg)$, higher-spin gravity is its value at $\cW_N$, and abelian
lattice holography is its class-G value. Drinfeld-Sokolov
functoriality links the four strata through the DS-Hochschild
chain-level bridge of Chapter~\ref*{ch:chiral-higher-deligne}.
